% answers_short.tex —— 2010 年·普通高等学校招生考试·重庆卷·理科 · 简短答案速查
% 由 tools/make_short_answers.py 自 answers.tex 的 \gksolhead/\gkanswer 提取，
% 并与 2010/2010_数学真题_重庆理_答案要点.md 逐题比对后生成，勿手改（改源 MD 后重跑）。

\begin{gknotice}
\textbf{使用说明}：本节只给最终答案便于快速核对，完整推导见解析版。答案为编写者独立推导，\textbf{非官方标准答案}。
\end{gknotice}

\gkpart{选择题}
\begin{center}\begin{tabular}{ccccccccccc}
\toprule
题号 & 1 & 2 & 3 & 4 & 5 & 6 & 7 & 8 & 9 & 10 \\
\midrule
答案 & A & B & B & C & D & D & B & C & C & D \\
\bottomrule
\end{tabular}\end{center}

\gkpart{填空题}
\begin{itemize}[leftmargin=2.6em]
  \item[\textbf{11.}] $-2\mathrm{i}$。
  \item[\textbf{12.}] $-3$。
  \item[\textbf{13.}] $\dfrac35$。
  \item[\textbf{14.}] $\dfrac83$。
  \item[\textbf{15.}] $\dfrac12$。
\end{itemize}

\gkpart{解答题}
\begin{itemize}[leftmargin=2.6em]
  \item[\textbf{16.}] （1）$[0,2]$；（2）$a=1$ 或 $a=2$。
  \item[\textbf{17.}] （1）$\dfrac45$；（2）$P(\xi=0)=\dfrac13$，$P(\xi=1)=\dfrac{4}{15}$，$P(\xi=2)=\dfrac15$，$P(\xi=3)=\dfrac{2}{15}$，$P(\xi=4)=\dfrac1{15}$，$E\xi=\dfrac43$。
  \item[\textbf{18.}] （1）$7x-4y-2=0$；（2）增区间 $(-1,1)$，$(7,+\infty)$；减区间 $(1,3)$，$(3,7)$。
  \item[\textbf{19.}] （1）$\sqrt3$；（2）$\dfrac{\sqrt6}{3}$。
  \item[\textbf{20.}] （1）$\dfrac{x^2}{4}-y^2=1$，渐近线 $y=\pm\dfrac12x$；（2）$2$。
  \item[\textbf{21.}] （1）$a_n=(n^2-1)c^n+c^{n-1}$；（2）$c\in\left(-\infty,-\dfrac{1+\sqrt{13}}{6}\right)\cup[1,+\infty)$。
\end{itemize}
